\documentclass{article}
\usepackage[margin=2cm]{geometry}
\usepackage{amsmath}
\usepackage{graphicx}
\usepackage{booktabs}
\usepackage[utf8]{inputenc}
\usepackage{ragged2e}
\setlength{\emergencystretch}{3em}
\begin{document}
\sloppy

\noindent\hfill\includegraphics[height=1.8em]{assets/figure_header_logo.png}\par\noindent\hrule\vspace{4pt}

$C . { \mathfrak { J } } ^ { ( \lambda t, L ; \iota ) } ( z )$ corresponding to a sphere with single puncture, by computing
residues.We define these to be


\begin{equation}
\begin{array} { r l } & { C _ { 3, \mathrm { s o n e m e n ~ }, \mathrm { ~ r o n e s e n c ~ e r e n e n ~ }, \mathrm { ~ r o n e s e n ~ }, \mathrm { ~ r o n e } } ^ { A ( L, L ) } } \\ & { C _ { 3, \mathrm { M } } ^ { B ( L, L ) } ( z ) = \underbrace { \left ( \prod _ { j \in 1 } \Gamma _ { c } \big ( a _ { i } / a _ { j } \big ) \right ) \Gamma _ { c } \big ( p q t ^ { 2}}}\end{array}
\end{equation}


The cap theory for zero values of $i f$ and $L$ is a model corresponding to
sphere with one puncture and flux$- \frac { 3 } { 4 }$ for U1)tfor$\frac { \underline { \mathfrak { d } } } { \mathfrak { u } }$ $\mathrm { U } ( 1 ) . { i }$and $- \frac { 1 } { 8 }$ for
$\mathrm { U } ( 1 ) . { j }$ See for details of the derivation of the flux. The index can be thought
as partition function on${ \mathfrak { S } } ^ { 1 } \times { \mathfrak { S } } ^ { 3 }$ and for non vanishing values of$i t$and the$L$
theory also has surface defects wrapping the${ \mathbb S } ^ { 1 }$ and one of the two equators.
of${ \mathfrak { S } } ^ { 3 }$ Finally we have a three punctured sphere $T . { \mathfrak { J } \mathfrak { J } \mathfrak { J } \mathfrak { C }. \mathfrak { J } \mathfrak { L } } ( \mathfrak { w }, \mathfrak { u }, \mathfrak { v } )$


\begin{equation}
\begin{array} { r l } & { \quad T _ { 3 a, \mathrm { J e }, \mathbf { J } _ { \rho } } ( w, u, v ) = \Gamma _ { c } \big ( \{ q p \} ^ { \frac { 1 } { 2 } t } \big \{ B ^ { - 1 } A \big ) ^ { \pm 1 } w ^ { \pm 1 } \big ) \Gamma _ { c } \left ( \frac { q p } { t ^ { 2 } } \right ) \big ( q ; q ) ( p ; p ) } \\ & { \times \int \! \! \! \frac { d h } { 4 \hbar i } \frac { \Gamma _ { c } \big ( p b ^ { \pm 1}}}}\end{array}
\end{equation}


where we have defined


\begin{equation}
\begin{array} { r l } & { H ( z _ { 1 }, z _ { 2 }, v _ { 1 }, v _ { 2 }, \alpha, b ; A ) = \{ g ; q \} ^ { 2 } \oint \frac { \mathrm { d } w _ { 1 } } { 4 i \pi w _ { 1 } } \oint \frac { \mathrm { d } w _ { 2 } } { 4 \pi w _ { 0 } } \int \frac { \mathrm { d } w _ { 2 } } { 2 \pi w _ { 0 } } \Gamma _ { \epsilon } \big ( \frac { p q t ^ { \frac { 3 } { 2 } } } { a ^ { + 1 } } w _}\end{array}
\end{equation}


The above expressions are non trivial to derive. The theory corresponding
to three punctured spheres is constructed by starting from a gauge theory,
index of which is roughly speaking   . and arguing that at some point on$H$
the conformal manifold the U(1) symmetry corresponding to fugacity $\sqrt { \mathfrak { a } / \mathfrak { b } }$
enhances to SU(2). This is a non trivial fact which follows from dualities.
This SU(2) is then taken to be dynamical with addition of some chiral fields.
The resulting index is given above. The statement that this theory corresponds to three punctured sphere is made by performing a variety of physical
consistency checks. Note that the construction also gives a theory having only
rank five symmetry as opposed to rank eight. For the three punctured sphere
we have fluxfor$3 / 4$ $\mathrm { U } ( 1 ) . { t }$ and vanishing flux for the Cartan generators of
SU(8). The three punctured sphere depends on four parameters $\{ 4, B, C ; D \}$
which parametrize SO(8) inside SU(8). That is,


\begin{equation}
\begin{array} { c c } {. } \\ { ( a _ { 1 }, a _ { 2 }, a _ { 3 }, a _ { 4 } ) = A ^ { \pm 1 } B ^ { \pm 1 }, \qquad ( a _ { 5 }, a _ { 6 }, a _ { 7 }, a _ { 8 } ) = C ^ { \pm 1 } D ^ { \pm 1 }. \qquad \quad ( 2 ) } \end{array}
\end{equation}

\end{document}
